\chapter{Wayland clipboard bridge}
\label{ch:wayland-clipboard}
\minitoc\label{wayland-clipboard-bridge}
\index{Wayland}\index{clipboard}\index{wl-copy}\index{interprogram-cut-function}\index{interprogram-paste-function}

% ── Opening signpost ──────────────────────────────────────────────
The previous chapters in this part dealt with universal Emacs
mechanisms.  This one addresses a platform-specific gap: terminal
Emacs on a Wayland session cannot reach the system clipboard through
X11 APIs.  This chapter describes the
\texttt{wl-copy}/\texttt{wl-paste} bridge and the two
\texttt{interprogram-*-function} hooks that make it transparent.

\section{The problem}\label{the-problem-2}

Emacs uses X11 clipboard APIs (\texttt{x-select-text}, \texttt{x-get-selection}) for system clipboard integration. When Emacs runs as a terminal client (\texttt{window-system} is \texttt{nil}) on a Wayland-native session (\texttt{WAYLAND\_DISPLAY} is set, no Xwayland), these X11 calls fail silently and the system clipboard is unreachable.

PGTK GUI Emacs (compiled with \texttt{-\/-with-pgtk}) has native Wayland clipboard support. The bridge below applies \textbf{only} to terminal Emacs (\texttt{emacs -nw}) on Wayland.

\section{The bridge implementation}\label{the-bridge-implementation}

\begin{minted}{elisp}
(when (and (eq system-type 'gnu/linux)
           (getenv "WAYLAND_DISPLAY")
           (not (display-graphic-p)))

  (defun my/wl-copy (text)
    "Send TEXT to the Wayland clipboard via wl-copy."
    (let ((proc (make-process :name "wl-copy"
                              :buffer nil
                              :command '("wl-copy" "-f" "-n")
                              :connection-type 'pipe)))
      (process-send-string proc text)
      (process-send-eof proc)))

  (defun my/wl-paste ()
    "Read the Wayland clipboard contents via wl-paste."
    (shell-command-to-string "wl-paste --no-newline"))

  (setq interprogram-cut-function   #'my/wl-copy
        interprogram-paste-function #'my/wl-paste))
\end{minted}

\section{\texttt{interprogram-cut-function} and \texttt{interprogram-paste-function}}\label{interprogram-cut-function-and-interprogram-paste-function}

These two variables are the system clipboard hooks:

\begin{itemize}
\item \texttt{interprogram-cut-function} is called with a string whenever text is added to the kill ring. Setting it to \texttt{my/wl-copy} propagates every kill to the Wayland clipboard.
\item \texttt{interprogram-paste-function} is called (no args) before \texttt{yank} to retrieve external clipboard content. Its return value, if non-nil and different from the top of the kill ring, is pre-pended to the ring and yanked.
\end{itemize}

\section{\texttt{make-process} with \texttt{:connection-type 'pipe}}\label{make-process-with-connection-type-pipe}

\texttt{make-process} spawns an asynchronous subprocess. The \texttt{:connection-type 'pipe} option connects the process's stdin to a pipe that Emacs can write to via \texttt{process-send-string}. The \texttt{-f} flag to \texttt{wl-copy} ("run in foreground") and \texttt{-n} ("no trailing newline") ensure the process stays alive exactly long enough to receive the text. \texttt{process-send-eof} closes the write end of the pipe, signalling end-of-input to \texttt{wl-copy}.

% ── Mental model summary ──────────────────────────────────────────
\section{Mental model}
PGTK GUI Emacs already has native Wayland clipboard support, so this bridge
is needed only for terminal Emacs (\texttt{emacs -nw}).  The integration
relies on two hooks: \texttt{interprogram-cut-function}, called on every
kill, is set to a function that pipes the killed text to \texttt{wl-copy};
\texttt{interprogram-paste-function}, called before every yank, is set to a
function that reads the current clipboard contents from \texttt{wl-paste}.
Under the hood, \texttt{make-process} with \texttt{:connection-type 'pipe}
connects the subprocess's stdin to a pipe that Emacs can write to via
\texttt{process-send-string}, and \texttt{process-send-eof} closes the write
end to signal end-of-input.
